\documentclass[12pt]{article}
\usepackage{geometry}
\usepackage{setspace}
\usepackage{titlesec}

\geometry{margin=1in}
\setstretch{1.3}

\begin{document}

\begin{center}

\vspace*{1.5cm}

{\LARGE \textbf{Design and Analysis of Sequential Circuit using Verilog}}\\[1.5cm]

{\large \textbf{Submitted By}}\\
Nirupam Das (2024CSB108)\\[0.8cm]

{\large \textbf{Guided By}}\\
Prof. Biplab Kumar Sikdar\\[2cm]

\end{center}

\section*{Introduction}

Sequential circuits are fundamental components of digital systems, as they store and process data based on both current inputs and previous states. Unlike combinational circuits, their behavior depends on clock signals and memory elements such as flip-flops.

The primary focus of this mini project is to design and analyze sequential circuits using Verilog Hardware Description Language (HDL), with a strong emphasis on leveraging Artificial Intelligence (AI) as a supportive tool in the hardware design process. In modern engineering workflows, AI significantly enhances learning, debugging, and optimization.

As part of this project, different sequential circuits were distributed among students to collectively explore a wide range of designs. My contribution focuses on the implementation of two important register-based circuits: the Universal Shift Register and the Bidirectional Shift Register.

The Universal Shift Register is capable of performing multiple operations such as shift left, shift right, parallel loading, and hold, making it highly versatile. The Bidirectional Shift Register enables shifting of data in both directions, which is essential in applications requiring controlled data movement.

Throughout the project, AI tools were utilized to improve conceptual understanding, assist in debugging Verilog code, and validate simulation outputs. This integration of AI enhanced efficiency and provided deeper insight into circuit behavior and design accuracy.

\section*{Tools and Technologies Used}

\begin{itemize}
    \item \textbf{Verilog HDL:} Used for designing and modeling sequential circuits.
    
    \item \textbf{Icarus Verilog (iverilog):} Used for compiling and simulating Verilog programs.
    
    \item \textbf{GTKWave / EPWave:} Used for waveform visualization and timing diagram analysis.
    
    \item \textbf{Text Editor / IDE:} Used for writing and editing code (e.g., VS Code).
    
    \item \textbf{Artificial Intelligence Tools:} Used for learning concepts, debugging errors, and improving design approach.
\end{itemize}

\section*{Methodology}

The methodology followed in this project combines structured hardware design principles with AI-assisted development techniques.

\begin{enumerate}
    \item \textbf{Understanding the Problem:}  
    The requirements of sequential circuits were studied, focusing on shift register operations and behavior.

    \item \textbf{Design Planning:}  
    The architecture of the Universal Shift Register and Bidirectional Shift Register was planned, including inputs, outputs, and control signals.

    \item \textbf{Verilog Implementation:}  
    The circuits were implemented using Verilog HDL, utilizing clock-driven sequential logic and case-based control structures.

    \item \textbf{AI-Assisted Development:}  
    AI tools were used to:
    \begin{itemize}
        \item Clarify theoretical concepts
        \item Debug logical and syntax errors
        \item Optimize code structure
        \item Analyze simulation results
    \end{itemize}

    \item \textbf{Simulation and Verification:}  
    Testbenches were created to verify functionality under different modes. Waveforms were analyzed to ensure correct operation.

    \item \textbf{Analysis:}  
    Outputs were compared with expected results and timing diagrams were examined to validate synchronous behavior.

\end{enumerate}

\end{document}